This appendix provides references to the key code artifacts produced in this project.

\section{Repository Links}

\begin{itemize}
    \item \textbf{RAG-Client-Knoccs:} The main repository containing the Frontend for the RAG system. \url{https://github.com/ibad23/RAG-Client-Knoccs} 
    \item \textbf{RAG-Server-Knoccs:} The main repository containing the RAG server, retrieval pipeline, and evaluation scripts. \url{https://github.com/fakehafaisal/RAG-Server-Knoccs}
    \item The actual Knoccs codebase is private and cannot be shared as it is owned by SpurSol.
\end{itemize}

\section{Sync Pipeline Trigger Script}

The sync pipeline is initiated by the following command, which reads the component mapping and generates a prompt for Claude Code:

\begin{lstlisting}[showstringspaces=false, frame=single, basicstyle=\ttfamily\scriptsize, caption={Sync Usage}, label={lst:sync}]
# Sync a specific component
npm run sync:figma toggle

# Generated files:
#   design-to-code/scripts/outputs/.figma-sync/sync-task.json
#   design-to-code/scripts/outputs/.figma-sync/prompt.txt

# Then run Claude with the generated prompt
claude "Read design-to-code/scripts/outputs/.figma-sync/prompt.txt
        and follow the instructions"
\end{lstlisting}

\section{RAG Retrieval Pipeline}

The following is a simplified representation of the hybrid retrieval logic from \texttt{search.py}:

\begin{lstlisting}[language=python, showstringspaces=false, frame=single, basicstyle=\ttfamily\scriptsize, caption={Hybrid Retrieval Pipeline (simplified)}, label={lst:rag_search}]
def search_and_summarize(self, query, top_k=15, initial_k=50,
                         conversation_history=None):
    # Step 1: Identify target company
    company_id = self._extract_company(query, conversation_history)
    if not company_id:
        return "Please specify which company you are asking about."

    # Step 2: Dual retrieval (BM25 + Dense)
    bm25_results = self.vectorstore.search_bm25(
        query, company_id, top_k=initial_k // 2)
    dense_results = self.vectorstore.search_dense(
        query, company_id, top_k=initial_k // 2)

    # Step 3: Merge and deduplicate
    combined = self._merge_and_deduplicate(bm25_results, dense_results)

    # Step 4: Cross-encoder re-ranking
    pairs = [(query, r["text"]) for r in combined]
    scores = self.reranker.predict(pairs)
    for r, s in zip(combined, scores):
        r["rerank_score"] = float(s)
    ranked = sorted(combined, key=lambda x: x["rerank_score"],
                    reverse=True)[:top_k]

    # Step 5: Assemble context and generate response
    context = "\n\n".join([r["text"] for r in ranked])
    prompt = self._build_prompt(query, context, conversation_history)
    response = self.llm.invoke(prompt)
    return response.content
\end{lstlisting}

\section{DeepEval Evaluation Configuration}

\begin{lstlisting}[language=python, showstringspaces=false, frame=single, basicstyle=\ttfamily\scriptsize, caption={DeepEval Metric Configuration (simplified)}, label={lst:deepeval}]
from deepeval.metrics import (
    AnswerRelevancyMetric, FaithfulnessMetric,
    ContextualPrecisionMetric
)
from deepeval.test_case import LLMTestCase
from deepeval import evaluate

threshold = 0.7
metrics = [
    AnswerRelevancyMetric(threshold=threshold, model="gpt-4o-mini"),
    FaithfulnessMetric(threshold=threshold, model="gpt-4o-mini"),
    ContextualPrecisionMetric(threshold=threshold, model="gpt-4o-mini"),
]

test_case = LLMTestCase(
    input=query,
    actual_output=rag_answer,
    retrieval_context=[chunk["text"] for chunk in retrieved_chunks],
    expected_output=ground_truth_answer
)

results = evaluate(test_cases=[test_case], metrics=metrics)
\end{lstlisting}

%%% Local Variables:
%%% mode: latex
%%% TeX-master: "../report"
%%% End:
